\documentclass[11pt,oneside]{book}
\usepackage{underscore}
\usepackage[margin=1in]{geometry}
\usepackage{xcolor}
\usepackage[hidelinks]{hyperref}
\usepackage{graphicx}
\usepackage{longtable}
\usepackage{booktabs}
\ifdefined\darkmode\pagecolor[HTML]{091117}\color[HTML]{F0EDE4}\def\viewtheme{dark}\else\def\viewtheme{light}\fi
\graphicspath{{../renders/}}
\newcommand{\cmd}[1]{\texttt{#1}}
\newcommand{\fig}[2]{\IfFileExists{../renders/#1}{\includegraphics[#2]{#1}}{\fbox{\parbox{0.8\linewidth}{Run \texttt{make views} to derive \path{renders/#1}.}}}}
\title{Maqam\\[0.5em]\large A live-coding terminal sequencer for just-intonation maqam,\\one specification and two implementations}
\author{Project team}
\date{\today}
\begin{document}
\frontmatter
\maketitle
\tableofcontents
\mainmatter

% ---------------------------------------------------------------------------
\chapter{What Maqam is}
maqam-live is a real-time terminal instrument for Arabic maqam music tuned in
just intonation. A performer types short phrases into a command box: a root, a
\emph{jins} (a scale fragment such as bayati or hijaz) and a rhythm written as
digits. Each phrase becomes a row of a timeline that the audio thread plays in
a loop. Control rows change tempo, sustain, filters and effects in no time;
jump rows repeat sections a counted number of times. The same timeline is saved
as a text session file, rendered offline to an MP4, edited by an LLM under a
checked allowlist, and run against a live guitar or microphone input through a
Neural Amp Modeler capture and a bank of sympathetic strings tuned by the score.

This project carries the upstream program (\path{rfielding/maqam}, revision
\texttt{c80b35d}) as a LeanFM project. Its record is this book together with
\path{spec/Requirements.lean} (the observable requirements, reverse-engineered
from the upstream source and documents), \path{spec/Requirements.proto} (the
typed messages), \path{spec/Maqam.leanfm} (the session type of the running
system) and one implementation mapping per implementation:

\begin{itemize}
  \item \path{impl/zig/} — a port to Zig, built, tested and run without the
  Rust: the default implementation and the reference for further ports;
  \item \path{impl/rust/} — the upstream Rust, the reference the requirements
  were distilled from.
\end{itemize}

Both are held to the same requirements. Review and acceptance happen in git;
chat histories are disposable.

\section{Running it}
From the project root:
\begin{itemize}
  \item \cmd{make spec-check} compiles the requirements and every
  implementation mapping in LeanFM's sandbox, validates each mapping against the
  one specification and regenerates \path{spec/tasks.json}.
  \item \cmd{make zig-check} builds and tests the default (Zig) implementation;
  \cmd{make app-check} builds and tests the Rust one. Neither touches the other's
  directory.
  \item \cmd{make compare} measures every implementation the same way
  (Chapter~\ref{ch:compare}).
  \item \cmd{make views} derives the figures of this book from the
  specification; \cmd{make book} rebuilds it in light and dark copies.
\end{itemize}

Either implementation runs as a terminal app with no arguments, or runs the
commands given as arguments (separated by \cmd{--}) without opening the
terminal UI, for example \cmd{maqam-live "load default.mq" -- "m 1"}.

% ---------------------------------------------------------------------------
\chapter{The music}
\section{Jins, maqam and just intonation}
A jins is a short scale fragment given as exact frequency ratios $p/q$ from its
root. The built-in ajnas (nahawand, bayati, hijaz, rast, kurd, saba, ajam,
nikriz, suznak, jiharkah, zamzam, zabal and others) can be abbreviated to an
unambiguous prefix; \cmd{create <Name> <p/q>...} adds a custom jins, which is
saved with the session. A comma stacks ajnas into one scale:
\cmd{d bayati, a nahawand 332}.

Roots are the letters \cmd{c d e f g a b}, with \cmd{+} or \cmd{-} for sharper
or flatter oud-lattice positions. Every root is snapped to a D-based oud
lattice (D4 = 293.6648\,Hz), and every phrase frequency is that root times one
of its jins ratios: the property \emph{Just intonation}.
\cmd{tuneto <pitch>} rebases the lattice so that pitch sounds at its 12-TET
frequency (A4 = 440\,Hz); it is live state, kept in \path{.globals.ml}, never in
a session.

The named church modes are Pythagorean, built from fifths alone: Minor and
Dorian are $1/1\;9/8\;32/27\;4/3\;3/2\;27/16\;16/9$, Ionian is
$1/1\;9/8\;81/64\;4/3\;3/2\;27/16\;243/128$, and Phrygian, Lydian, Mixolydian,
Aeolian and Locrian use the same 3-limit ratios (256/243, 128/81, 729/512,
1024/729). Upstream used 5-limit ratios (6/5, 8/5, 9/5) for these until
revision \texttt{a2cdac8}, so \cmd{d minor} sounds different in sessions made
before it.

\section{Rhythm and timing}
A rhythm is a string of non-zero digits, each a group size. The first
subdivision of a group is a kick, the rest snares: \cmd{332} is
\texttt{X..X..X.}. One digit unit is an eighth note, $60/(2\cdot\mathit{bpm})$
seconds. A phrase without a rhythm inherits the last one. \cmd{r<N>} repeats a
phrase $N$ times (at most 20).

\section{The signal chain}
In each audio callback the voices of the playing phrases, the sub-bass and the
drums are mixed with the live input, the sympathetic strings excited by that
input and tuned to the current phrase, the per-instrument VCF bank and the
effects (flanger, chorus, reverb, delay, ping-pong), then soft-clipped and
clamped to $[-1,1]$: the required proof \emph{Output is bounded}. Each
implementation's mapping names the code of this mix loop.

\section{The reverb}\label{sec:reverb}
The reverb is the specification's own, not upstream's: upstream's was two
feedback echoes (29 and 41\,ms, two of its four taps) writing into one shared
buffer, clamped to $[-1,1]$ inside the loop, with no damping or diffusion, so it
rang at about 34 and 24\,Hz and distorted loud tails. Both implementations,
the Rust snapshot included, compute the same replacement, so renders agree. Its commands, ranges
and defaults are upstream's (\cmd{reverb mix 0..1 decay 0..0.98}). Decay $d$
sets the time the tail takes to fall 60\,dB, $T_{60} = 0.25 \cdot 16^{d/0.98}$\,s:
0.25\,s at 0, 1.6\,s at the default 0.65, 4\,s at 0.98. Each channel first
passes two one-pole high-pass stages at 120\,Hz, so a bass's fundamentals stay
dry and the tail sits above them, then four allpass diffusers (4.77, 3.59,
12.73, 9.31\,ms; the right channel's 1.07 times longer), which smear an attack
into a dense wash from the first milliseconds. The tank is eight delay lines
(29.7 to 73.3\,ms) mixed by an orthogonal Householder matrix; each line loses
the gain that makes it fall 60\,dB in $T_{60}$ and is damped by a one-pole
low-pass at 5\,kHz, so the highs fade before the lows, as in a room. The left
input feeds the even lines and the right the odd lines; the left output sums
the even lines and the right the odd ones, which leaves them uncorrelated and
wide. The wet sum is scaled by $0.6/\sqrt{T_{60}}$ so \cmd{mix} means the same
at every decay. Nothing inside the reverb clamps: the only clamp is the final
one on the output (\emph{The reverb never distorts or rings}), and state that
fades below $10^{-15}$ is set to zero, so a fading tail never slows a block
with subnormal floats. The full definition is the markdown section
\texttt{reverb} of \path{spec/Requirements.lean}.

The live input takes one path. With NAM on and a capture loaded only the
modeled signal goes on, with no clean signal beside it; with NAM off the clean
input goes on instead. Next comes the cabinet stage (\cmd{cab vintage},
\cmd{cab modern} or \cmd{cab open}, with an optional \cmd{mix}), a short
impulse response that is bypassed while off. Without \cmd{vcf all} the result
is the mic bus. Every bus is scaled by its own volume after its filter
(\cmd{vol mic 0.25}, \cmd{vol bass 2}); the mic bus is added after the soft
clip, so its trim is linear. Last, \cmd{dry} $d$ crossfades the whole
output with the raw input, untouched by NAM, filters, effects or volume:
$\mathit{out} = \mathit{processed}\cdot(1-d/100) + \mathit{raw}\cdot d/100$.
Captured frames never queue up behind more than three output blocks
(\emph{Live input never lags behind a backlog}), and \cmd{--bufsize=<frames>}
asks the device for its buffer size (default 64).

% ---------------------------------------------------------------------------
\chapter{The command language}\label{ch:commands}
A line is split on \cmd{;} (no escaping) and each part is one command. A parse
or range error is shown and stops the rest of the line; nothing it would have
changed is changed (\emph{A rejected command changes nothing}).
\cmd{clockin} and \cmd{clockout} are taken before the parser. A line whose text
before the first \cmd{:} is \cmd{chatgpt}, \cmd{gpt}, \cmd{openai},
\cmd{claude} or \cmd{anthropic} is a prompt for an LLM (Chapter~\ref{ch:tasks}).

\begin{longtable}{p{0.36\linewidth}p{0.58\linewidth}}
\toprule
Command & Meaning\\
\midrule
\cmd{<root> <jins> [rhythm] [, <root> <jins>] [rN]} & a phrase row\\
\cmd{j <id|start> [times|always]} & jump back to a row; \cmd{times} passes, then fall through\\
\cmd{j <ok> else <fail> [times]} & branch on counter success or exhaustion\\
\cmd{i <id> <cmd>}, \cmd{edit <id> <cmd>} & insert before a row, replace a row\\
\cmd{x <id>...}, \cmd{up}/\cmd{down <id>}, \cmd{rot} & delete, move, rotate rows (ids stay)\\
\cmd{z [id|start]}, \cmd{start}, \cmd{pause}, \cmd{stop} & playback and the stop row\\
\cmd{bpm}, \cmd{s}/\cmd{sus} & tempo and sustain rows\\
\cmd{vol [<bus>] <n>}, \cmd{dry <\%>}, \cmd{tuneto} & live volume (overall or per bus: all, mic, bass, kanun, drums, sym), dry mix and tuning, kept in \path{.globals.ml}; inserted with \cmd{i} or \cmd{edit}, vol and dry are timeline rows\\
\cmd{vcf}/\cmd{cut}/\cmd{res}/\cmd{drive} & per-instrument filter rows\\
\cmd{flanger}, \cmd{chorus}, \cmd{reverb}, \cmd{delay}, \cmd{pingpong}, \cmd{cab}, \cmd{fx off} & effect and cabinet rows\\
\cmd{sym ...} & sympathetic strings (harmony, decay, gain, per-source partitions)\\
\cmd{nam ...} & NAM captures: import, pin, search, gain, input, off\\
\cmd{save [path]}, \cmd{load <path>}, \cmd{ls} & sessions\\
\cmd{m[N]} & record $N$ cycles to an MP4\\
\cmd{create <Name> <p/q>...}, \cmd{delete <Name>} & custom ajnas\\
\cmd{audition}, \cmd{clear}, \cmd{?}/\cmd{help}, \cmd{q}/\cmd{quit} & the rest\\
\bottomrule
\end{longtable}

Numeric values may be absolute, \cmd{+n}, \cmd{-n}, \cmd{*k} or \cmd{/k}, or a
per-tick step \cmd{nt}; \cmd{+0} stops a tick. Out-of-range values are rejected
with a message naming the range:

\begin{longtable}{p{0.36\linewidth}p{0.58\linewidth}}
\toprule
Setting & Range\\
\midrule
bpm & 20..400 (default 120)\\
sustain & 0.05..10\,s (default 1.25)\\
vol, vol per bus & 0..2 (default 1)\\
dry & 0..100\,\% (default 0)\\
cab mix & 0..1 (default 1)\\
vcf cutoff, resonance, drive & 10..22000\,Hz (at most 0.45 of the sample rate), 0..0.98, 0.1..12\\
flanger rate, depth, delay, feedback, mix & 0.01..8\,Hz, 0..1, 0.1..10\,ms, $-0.95$..0.95, 0..1\\
chorus delay; reverb decay & 5..35\,ms; 0..0.98\\
delay time, feedback & 0.01..2\,s, 0..0.95\\
sym decay, gain, interval ratio & 0.9..0.99999 per ms (default 0.999), 0..512 (default 2), 1/4..4, at most 8 intervals\\
nam gain & 0..100\\
\bottomrule
\end{longtable}

Entry ids are stable: insert, delete, move and rotate never renumber an
existing id, and jumps name ids, not positions (\emph{Entry ids are stable}).

% ---------------------------------------------------------------------------
\chapter{Sessions}
\cmd{save} writes a \texttt{MAQAM\_SESSION\_V3} file: a header line, a
\cmd{create} line per custom jins, then one record per timeline row in order:

\begin{center}\begin{tabular}{ll}
\toprule
Record & Row\\
\midrule
\texttt{P|id|repeat|src} & phrase\\
\texttt{J|id|target|times}, \texttt{J|id|target|fail|times} & jump\\
\texttt{B|id|bpm}, \texttt{S|id|sustain} & tempo, sustain\\
\texttt{V|id|vcf}, \texttt{F|id|fx}, \texttt{N|id|nam}, \texttt{Y|id|sym} & filter, effects, NAM, sympathetic strings\\
\texttt{L|id|vol}, \texttt{D|id|dry} & volume (overall or per bus), dry mix\\
\texttt{T|id|stop} & stop\\
\bottomrule
\end{tabular}\end{center}

\verb!\!, \verb!|! and newline are escaped. V1 and V2 files still load, with
sequential ids. \cmd{load} rebuilds the audio from the timeline (Clear, the
start settings, then every entry) and reports a failure with its line number.
\emph{Session round trip}: loading a saved timeline gives the same entries,
ids, sources and custom ajnas; \cmd{vol} and \cmd{tuneto} are never in a
session.

% ---------------------------------------------------------------------------
\chapter{How the system works}\label{ch:tasks}
Eight actors take part: the \emph{Performer} at the terminal, the \emph{App}
(the timeline, the command dispatch and the workers), the \emph{Audio} output
callback, the \emph{Recorder} and the \emph{Ffmpeg} process it starts, an
\emph{Llm} provider, the \emph{Web} (NAM captures, search, TONE3000) and the
\emph{Disk}. The App owns the timeline; the Audio thread only receives
commands over a bounded channel (512 commands) and publishes its position
through atomics, so the callback never blocks (\emph{The audio callback never
blocks}).

\section{The whole protocol}
Figure~\ref{fig:protocol} is LeanFM's sequence diagram of protocol
\texttt{MaqamLive} in \path{spec/Maqam.leanfm}, generated by
\texttt{session\_diagram}; LeanFM's session view animates it and checks event
logs against it.

\begin{figure}[htbp]\centering
\fig{maqam-live-\viewtheme.pdf}{height=0.82\textheight,keepaspectratio}
\caption{Protocol MaqamLive (generated).}\label{fig:protocol}
\end{figure}

\section{The tasks}
Each task is a state machine derived from the requirements' grammars
(\path{spec/tasks.json}).

\begin{description}
  \item[live\_edit] A command line is parsed; a rejection is shown and changes
  nothing, otherwise the App changes its timeline first and then tells the
  audio (\emph{Timeline first, audio second}).
  \item[timeline\_playback] The callback walks the rows: a phrase plays its
  repeats, a control applies in no time (\emph{Controls take no time}), a jump
  row counts passes — \cmd{j X N} passes it $N$ times: $N-1$ jumps taken, then
  one exhausted that resets its counter; \cmd{always} never exhausts
  (\emph{Jump passes}).
  \item[session\_save\_load] Load reads the whole file, then either rebuilds the
  audio and reports the row count, or reports the failing line.
  \item[record\_mp4] The timeline is rendered offline, encoded by ffmpeg
  (\texttt{nice -n 10}, one retry without subtitles) into a staged
  \path{.partial.mp4} that is renamed to \path{./maqam-<secs>.mp4}; a failed
  recording leaves no partial file. ffmpeg missing from PATH is reported with
  what to install.
  \item[llm\_edit] \cmd{chatgpt:} or \cmd{claude:} asks the provider with the
  language reference and the score as context. Returned tool commands are
  applied only when every one parses and is on the allowlist; otherwise none is,
  and any failure while applying restores the previous score
  (\emph{LLM commands are checked before any is applied}). The LLM can never
  save, load, record, play, quit or change \cmd{vol} or \cmd{tuneto}.
  \item[nam\_fetch] A capture is downloaded into \path{<name>.nam.part},
  renamed into the cache (\path{./.nam} or \texttt{MAQAM\_NAM\_CACHE\_DIR})
  when complete, and only then loaded onto the live input (\emph{A download is
  cached before it is loaded}); cache names are sanitized so no write leaves
  the cache.
\end{description}

\begin{figure}[htbp]\centering
\fig{live_edit.pdf}{width=0.45\linewidth}\hfill\fig{timeline_playback.pdf}{width=0.45\linewidth}\\[1em]
\fig{session_save_load.pdf}{width=0.45\linewidth}\hfill\fig{record_mp4.pdf}{width=0.45\linewidth}\\[1em]
\fig{llm_edit.pdf}{width=0.45\linewidth}\hfill\fig{nam_fetch.pdf}{width=0.45\linewidth}
\caption{The six tasks' state machines, derived from \texttt{spec/tasks.json}.}
\end{figure}

% ---------------------------------------------------------------------------
\chapter{Requirements}
\section{Positive invariants}
Always: a rejected command changes nothing; the timeline changes before the
audio is told; entry ids are stable; jump passes count as above; controls take
no time; every frequency is just; a session round-trips; LLM commands are all
checked before any is applied; a download is cached before it is loaded; with
\texttt{MAQAM\_EVENT\_LOG} set every message of \texttt{MaqamLive} the app takes
part in is logged once, in order, conforming to the protocol; the offline render
is deterministic; every audio block is on time (Section~\ref{sec:deadline});
the reverb's tail falls 60\,dB in its $T_{60}$, is dense from 100\,ms on and
is uncorrelated between left and right (Section~\ref{sec:reverb}).

\section{Negative requirements}
Never: the LLM saves, loads, records or plays; a failed recording leaves a
partial MP4; the audio callback blocks; the audio thread waits on the kernel (a
page fault in the callback, or normal priority where real-time is allowed); a
buffer skip that is not counted and shown; an output sample outside $[-1,1]$; an
input that crashes the app; a secret (API key, TONE3000 token) in a session, the
terminal, the LLM context, a recording or a log; a cache write outside the
cache; tests that open a browser or reach the network; a skipped block at 64
frames and 48\,kHz with NAM, the strings and the reverb on; an event log
field the protocol does not declare; a reverb that clips or rings inside its
loop, muddies a 40\,Hz note (less than 12\,dB below a 1\,kHz one) or leaves
subnormal floats in its state; and one implementation that reads, builds
or links another's directory or needs another language's toolchain
(\emph{Implementations stand alone}).

\section{The audio deadline is hard}\label{sec:deadline}
\emph{Eventually} is not enough for audio. A callback that finishes 10\,ms late
satisfies ``eventually the block is ready'' and is still a gap the listener
hears. The requirement is a bounded one, and it repeats forever, on the
device's clock: the device releases block $k$ at $t_0 + k\,P$ with
$P = 64/48000$\,s $= 1333.3\,\mu$s, whether the audio is ready or not, and
plays it only if it was handed over before $t_0 + (k+1)\,P$; otherwise it skips
it. A bounded deadline is a safety property: a miss is seen at a finite time,
so it can be checked block by block from a log, not only argued.

Protocol \texttt{AudioPeriod} (\path{spec/AudioPeriod.leanfm}) states one
period: \texttt{Pull} from the device, \texttt{Block} from the audio, then the
device decides \texttt{Played} or \texttt{Skipped}. Each block is its own
session, so a run of the bench is a stream of independent trials and one
\texttt{Skipped} anywhere breaks \emph{No missed audio deadline}. The verdict
is exact, in integers: $(\mathit{Block}-\mathit{Pull}) \cdot 48000 \le 64
\cdot 10^9$\,ns. Every block counts: the maximum, never a percentile.

On time means woken on time as well as rendered in time, and that puts two
more requirements on the audio thread: it never takes a page fault in the
callback, and it runs with real-time priority where the machine allows it
(\emph{The audio thread never waits on the kernel}). The Zig bench found both:
at normal priority one release was 3.2\,ms late, and the first block after a
load took 916 page faults (1\,ms), each a skip. With memory locked and
real-time scheduling (RealtimeKit), three runs of 7500 blocks had none late.
macOS has no \texttt{mlockall}; there the buffers the callback uses are written
once before the device starts, and real-time scheduling is the Mach
time-constraint policy (Chapter~\ref{ch:compare}).
On the live device a skip is never silent: each underrun is counted and shown
in red in the header.

\section{Hooks every implementation provides}
\begin{itemize}
  \item \textbf{Event log}: with \texttt{MAQAM\_EVENT\_LOG=<path>}, one JSON
  object per line for each message of \texttt{MaqamLive} the app takes part in:
  \texttt{\{ts, session, from, to, label, cause, fields\}}, \texttt{ts} in
  nanoseconds of one monotonic clock, \texttt{cause} the \texttt{ts} of the
  event the protocol names as its cause, only the declared fields, integers as
  decimal strings. Unset, nothing is written.
  \item \textbf{Offline render}: \cmd{maqam-live --render <session.mq> --cycles
  <n> --rate <hz> --wav <out.wav>} writes 32-bit float mono WAV without opening
  audio devices or the terminal UI, the same samples on every run.
  \item \textbf{Bench}: \cmd{maqam-live --bench <session.mq> --seconds <s>
  --frames 64 --rate 48000 [--nam <capture.nam>] [--period-log <path>]} runs
  the live audio path against a null device with the strings and the reverb on.
  The null device is the clock: it releases each block on the period grid, the
  bench waits for each release, and times the block from release to hand-over,
  with real-time scheduling and locked memory where allowed. It prints one JSON
  line \texttt{\{frames, rate, callbacks, deadline\_misses, callback\_p50\_us,
  callback\_p99\_us, callback\_max\_us, release\_late\_max\_us, slack\_min\_us,
  realtime, memory\_locked\}}, and with \cmd{--period-log} writes the
  \texttt{AudioPeriod} stream after the run.
\end{itemize}

\IfFileExists{implementations.tex}{% Included by main.tex. One specification, two implementations.
\chapter{Implementations and how they compare}\label{ch:compare}

The requirements in Chapter~\ref{ch:tasks} and the specification files fix what
Maqam does; they say nothing about the language. \path{manifest.json} lists two
implementations of the same specification, each mapped to it by its own
\path{impl/<name>/Implementation.lean} and checked by \cmd{make spec-check}.

\begin{description}
  \item[\path{impl/rust/source}] the upstream Rust, the reference the
  requirements were distilled from, with one change: its reverb is the
  specification's (Section~\ref{sec:reverb}). Built with \cmd{cargo}.
  \item[\path{impl/zig/source}] the default implementation, a port to Zig 0.17 that runs on Linux and
  macOS. On Linux it needs only the Zig compiler and the ALSA development
  headers (\cmd{libasound2-dev}); on macOS only the Zig compiler and the
  system's CoreAudio, with the Command Line Tools' SDK. ffmpeg is needed at run
  time for \cmd{m} (recording), as upstream.
\end{description}

\section{Building the Zig implementation on its own}
Nothing under \path{impl/rust/} is read, built or linked by the Zig build
(proof \emph{Implementations stand alone}). The bundled \path{.mq} sessions are
copied into \path{impl/zig/source} because they are inputs of the
specification, not of the Rust build. From \path{impl/zig/source}:

\begin{center}
\begin{tabular}{ll}
\toprule
\cmd{make} & build \path{zig-out/bin/maqam-live} (ReleaseFast)\\
\cmd{make check} & build and run every module's tests\\
\cmd{make render SESSION=default.mq} & offline render to \path{render.wav}\\
\cmd{make bench SESSION=default.mq NAM=x.nam} & the real-time bench at 64 frames\\
\cmd{make run} & the terminal app\\
\bottomrule
\end{tabular}
\end{center}

\section{One device per operating system}
The engine never knows which device drives it: \texttt{device.Output} is the
ALSA device on Linux and the CoreAudio device on macOS, chosen when the program
is compiled, and each calls \texttt{audio.Engine.process} once per block. They
differ in who keeps time. ALSA is written to: the device thread renders a block
and blocks in the write until the hardware has room. CoreAudio pulls: it calls
the output callback on its own real-time thread when it needs a block, and a
second, input-only unit hands the captured frames to that callback through a
lock-free ring. The requirements on the callback are the same on both.

The controls behind \emph{The audio thread never waits on the kernel} and
\emph{A buffer skip is never silent} are the operating system's:

\begin{center}
\begin{tabular}{p{0.24\linewidth}p{0.33\linewidth}p{0.33\linewidth}}
\toprule
 & Linux (ALSA) & macOS (CoreAudio)\\
\midrule
real-time scheduling & \texttt{SCHED\_FIFO}, else RealtimeKit (\texttt{fifo}, \texttt{rtkit}) & the Mach time-constraint policy (\texttt{mach}); CoreAudio's own I/O thread already has it\\
no page faults & \texttt{mlockall}, current and future mappings & no \texttt{mlockall} on macOS (\texttt{memory\_locked} false): the callback's buffers are written once before the device starts\\
a skip is counted & the write returns \texttt{-EPIPE} & the device's sample time jumps past the frames handed over\\
opening a URL & \cmd{xdg-open} & \cmd{open}\\
\bottomrule
\end{tabular}
\end{center}

\noindent On macOS the terminal running \cmd{maqam-live} needs microphone
permission for the live input; without it the app plays with no input, as on a
Linux machine with no capture device.

\noindent From the project root \cmd{make zig-check} does the same as
\cmd{make check} there, and \cmd{make check} runs the specification check and
both implementations' checks.

\section{How the Zig is arranged}
Each Rust module has a Zig module of the same name and role; types keep their
Rust names and functions become camelCase (\path{impl/zig/source/PORTING.md}).

\begin{longtable}{lll}
\toprule
Rust & Zig & role\\
\midrule
\endhead
\texttt{tuning.rs} & \texttt{tuning.zig} & pitches, ajnas, the 5-limit lattice\\
\texttt{command.rs} & \texttt{command.zig} & the command language and its messages\\
\texttt{sequencer.rs}, \texttt{synth.rs} & same names & phrases, bars, walks, voices\\
\texttt{session\_v3.rs} + loaders in \texttt{app.rs} & \texttt{session.zig} & MAQAM\_SESSION\_V3/V2/V1\\
\texttt{fx.rs}, \texttt{vcf.rs}, \texttt{analog.rs}, \texttt{sympathetics.rs} & same names & the signal chain\\
crate \texttt{nam-rs} 0.4 & \texttt{nam.zig} & NAM WaveNet (A1, A2), LSTM, slimmable\\
\texttt{audio.rs} + \texttt{cpal} & \texttt{audio.zig}, \texttt{device.zig}, \texttt{coreaudio.zig} & the engine, ALSA and CoreAudio\\
\texttt{record\_old.rs}, \texttt{carpet.rs} & \texttt{record.zig}, \texttt{carpet.zig}, \texttt{png.zig} & the MP4 recorder\\
\texttt{app.rs} & \texttt{app.zig}, \texttt{net.zig} & the App and its workers\\
\texttt{ui.rs} + \texttt{ratatui} & \texttt{ui.zig}, \texttt{term.zig} & the terminal UI\\
crossbeam channels, main.rs statics & \texttt{channel.zig}, \texttt{shared.zig} & lock-free rings, atomics\\
--- & \texttt{eventlog.zig}, \texttt{wav.zig} & the comparison hooks\\
\bottomrule
\end{longtable}

The audio thread in the Zig never allocates, frees or locks. The App sends
pointers to phrases and models on a 512-slot lock-free ring; the engine sends
what it no longer uses back on a second ring and the App frees it. Upstream's
mutex-guarded jump-counter map is a fixed table of atomics.

\section{Where the Zig differs on purpose}
Messages, ranges, session files and commands are the same. The differences are
the requirements the Rust does not meet (Chapter~\ref{ch:gaps}):

\begin{itemize}
  \item A phrase whose combined scale has fewer than three notes is rejected
  with a message instead of panicking (\cmd{create Zz 1/1}, \cmd{d zz 4}).
  \item Inputs that loop forever or crash upstream are handled: a non-ASCII
  value token, deeply nested \cmd{edit}/\cmd{insert}, an infinite or zero
  frequency in the tuning lattice or the sympathetic strings, and NAM files
  that make \texttt{nam-rs} index out of bounds.
  \item Jins name ties resolve alphabetically instead of by hash order, so
  \cmd{d bay} always means the same jins.
  \item Bounds: repeats at most \texttt{r20}, jump times at most 10000, the
  timeline at most 4096 rows, LLM responses at most 4\,MiB, NAM downloads at
  most 256\,MiB, NAM files at most 64\,MiB, LLM context at most 16
  exchanges, line history at most 1000 lines, recordings at most 600\,s, and
  the voice lists, jump counters and garbage queues fixed in advance.
  \item A recording that fails leaves no \path{.partial.mp4} and no temporary
  layers; after an ffmpeg failure only the \path{maqam-ffmpeg.log} the message
  points to is kept.
  \item Every NAM cache name is sanitized; tests never reach the network or
  open a browser.
  \item MIDI clock in and out (\cmd{clockin}, \cmd{clockout}), which the
  specification does not mention, answer that MIDI is not available.
  \item The audio thread locks its memory (Linux) and asks for real-time
  scheduling before its loop, and the live device counts every underrun and shows it in
  the header; upstream recovers from an underrun silently.
  \item One regression: Zig 0.17's HTTP client has no request timeouts, where
  upstream gives up after 30\,s (LLM) and 120\,s (downloads). Only the worker
  thread waits; the audio and the terminal stay live.
  \item The comparison hooks (\cmd{--render}, \cmd{--bench},
  \texttt{MAQAM\_EVENT\_LOG}) exist.
\end{itemize}

\section{Measuring them the same way}
\cmd{make compare} in the project root runs each implementation's
\cmd{check}, \cmd{render} (twice, to test determinism) and \cmd{bench}, checks
the event log against protocol \texttt{MaqamLive} with LeanFM's
\texttt{session\_conform}, checks every block of the bench's \texttt{AudioPeriod}
stream (exactly, and its conformance; Section~\ref{sec:deadline}), null-tests each render against the default
implementation's, and writes \path{reports/compare/}. A target an
implementation does not provide is reported as \emph{not provided}, never as a
pass; upstream Rust provides no hooks, so its render and bench columns say so.

\IfFileExists{../reports/compare/summary.tex}{% Generated by scripts/compare.py; do not edit.
\noindent Measured by \cmd{make compare}: session \path{impl/zig/source/default.mq}, 2 cycles at 48000\,Hz; bench 30\,s of 64-frame blocks at 48\,kHz with the sympathetic strings and the reverb, NAM \path{impl/zig/source/testdata/nam/marshall_jcm800.nam}.

\begin{longtable}{lp{0.78\linewidth}}
\toprule
\multicolumn{2}{l}{\textbf{zig (default)}}\\
specification & validated\\
check & pass 0.1 s\\
render & pass, 14.8 s, deterministic\\
event log & conforms\\
bench & 22500 blocks, 0 deadline misses; p50 218 us, p99 301 us, max 791 us; slack 542 us; scheduling rtkit, memory locked\\
every block & 0 of 22500 late, 0 labels disagree; AudioPeriod: conforms\\
source lines & 27572\\
\midrule
\bottomrule
\end{longtable}
}{%
\noindent Run \cmd{make compare} to fill this section with measurements.}
}{}

% ---------------------------------------------------------------------------
\chapter{Known gaps}\label{ch:gaps}
These are counterexamples to the requirements in the upstream Rust
(\path{impl/rust}); the requirements still forbid them, and the Zig
implementation avoids each.

\begin{description}
  \item[A one-note jins panics] \cmd{create Zz 1/1} then \cmd{d zz 4} panics in
  \texttt{build\_phrase} (\texttt{clamp} with min $>$ max)
  (\path{reports/reverse-engineering/panic-one-note-jins.txt}). Breaks
  \emph{No input crashes the app}.
  \item[An unsanitized cache name] A \cmd{nam tone3000} name read from a loaded
  \path{.mq} reaches the cache path without sanitizing. Breaks \emph{Cache
  writes stay in the cache}.
  \item[A test that reaches out] Loading the bundled \path{nama2.mq} in a test
  starts the TONE3000 login and runs \texttt{xdg-open}. Breaks \emph{Tests do
  not reach out}.
  \item[Reverb stutters at 64 frames] The reverb misses deadlines at 64-frame
  buffers. Breaks \emph{No missed audio deadline}.
  \item[A cabinet change allocates in the callback] Changing the cabinet model
  rebuilds its impulse response and resizes its history inside the output
  callback (\texttt{CabinetProcessor::set\_model}); the allocator may take a
  lock there. Breaks \emph{The audio callback never blocks}. The Zig builds all
  three responses before the device starts.
  \item[The tests deadlock] \texttt{bundled\_v3\_sessions\_load} gives the App a
  32-slot audio channel nobody reads; loading \path{default.mq} sends more
  commands than that (\texttt{SetDry} among them since revision
  \texttt{72c11f5}), so the send blocks forever while the test holds the
  session lock, and every test that takes the lock waits behind it:
  \cmd{make app-check} never finishes. \texttt{recording\_errors\_appear\_in\_response\_area}
  never finishes either, before and after the reverb change; with both skipped
  and one test at a time, the other 80 pass. Run in parallel, tests fail
  together behind the shared session lock
  (\path{reports/zig-macos/catch-up-72c11f5.txt}).
  \item[No comparison hooks] The upstream binary has no \cmd{--render},
  \cmd{--bench} or event log.
  \item[Unbounded] Timeline length, phrase repeats, jump times, the recording
  buffer, LLM history and NAM download size are not bounded upstream.
\end{description}

\noindent One gap is the Zig implementation's own, on macOS only:

\begin{description}
  \item[Late wake-ups on macOS] Under the Mach time-constraint policy the bench
  thread is still sometimes woken up to 0.5\,ms after its release on an Apple M1
  desktop, so a few of 7500 blocks are skipped (0 to 11 per run, with and without
  a NAM capture); rendering itself stays within the period
  (\path{reports/zig-macos/bench-m1.txt}). The live callback runs on
  CoreAudio's own I/O thread, which counted no skip in the same session. Breaks
  \emph{No missed audio deadline} for the bench on macOS.
\end{description}

\chapter{Where the documents and the code disagree}
Recording does not reuse the live path: it has its own mix loop at a fixed
44.1\,kHz without the sympathetic strings, NAM, microphone or \cmd{vol},
peak-normalizes to 0.9 and assumes 0.25\,s subdivisions when spawning voices.
NAM rows are saved (\texttt{N|}) although the upstream documents call NAM live
state. The upstream \path{docs/session-v3.md} lists only B, S, P, J, T. Upstream's
reverb, which uses two of its four taps, is not followed: the specification
replaces it (Section~\ref{sec:reverb}). The README writes Saba's endpoint as 80/64
where the code says 5/4, and the commit that added \cmd{--bufsize} says its
default is 256 where the code uses 64. Where code and documents disagree the requirements follow the
code, so every implementation behaves as the upstream code does.

\end{document}
